\documentclass[12pt]{article}

\usepackage[margin=1in]{geometry}
\usepackage{amsmath,amssymb}
\usepackage{xcolor}

% colours matching the slide
\definecolor{wcol}{HTML}{E0863B}   % weights  (warm)
\definecolor{dcol}{HTML}{3FB97E}   % data / keys / query (cool green)
\definecolor{vcol}{HTML}{3E7BD6}   % values / error signal (blue)

\newcommand{\W}[1]{\textcolor{wcol}{#1}}
\newcommand{\D}[1]{\textcolor{dcol}{#1}}
\newcommand{\V}[1]{\textcolor{vcol}{#1}}

\setlength{\parindent}{0pt}
\setlength{\parskip}{6pt}

\title{Slide 22: connecting line 1 to line 2}
\author{}
\date{}

\begin{document}
\maketitle

\textbf{The two lines on the slide.}

\medskip

Line 1 (self-attention):
\[
  F(\D{q}) \;=\; \W{W}_{\!V}\,\D{P}\,(\W{W}_{\!K}\,\D{P})^{\top}\,\D{q}
\]

Line 2 (after splitting the prompt):
\[
  F(\D{q}) \;=\; \W{W}_{\mathrm{ZS}}\,\D{q}
  \;+\; \sum_{i}(\W{W}_{\!V}\,\D{x_i})(\W{W}_{\!K}\,\D{x_i})^{\top}\,\D{q}
\]

The jump hides two moves: a matrix identity, then a split.

\hrulefill

\section*{The identity you're missing}

A matrix times another matrix's transpose equals the sum of the outer
products of their columns. For two $d\times N$ matrices
$A=[a_1\ \cdots\ a_N]$ and $B=[b_1\ \cdots\ b_N]$ (columns $a_j,b_j$):
\[
  A\,B^{\top} \;=\; \sum_{j=1}^{N} a_j\,b_j^{\top}.
\]

Check on a single entry $(r,c)$:
\[
  (A B^{\top})_{rc}
   = \sum_{j} A_{rj}\,B_{cj}
   = \sum_{j} (a_j)_r\,(b_j)_c
   = \sum_{j} (a_j b_j^{\top})_{rc}. \qquad\checkmark
\]

This is the same outer-product fact that built the trained weights on the
previous slide ($\W{W}_T=\W{W}_0+\sum_i \V{\sigma_i}\D{x_i}^{\top}$) --- here it
runs in reverse: a matrix product \emph{unpacks into} a sum of outer products.

\section*{Step 1 — regroup line 1 as ``build a matrix, then apply it to $\D q$''}

Matrix multiplication is associative, so bracket line 1 as the effective
matrix first:
\[
  F(\D q)=\Big[\underbrace{(\W{W}_{\!V}\D P)(\W{W}_{\!K}\D P)^{\top}}_{d\times d \text{ matrix}}\Big]\,\D q .
\]

Apply the identity with $A=\W{W}_{\!V}\D P$ (columns $\W{W}_{\!V}\D{p_j}$) and
$B=\W{W}_{\!K}\D P$ (columns $\W{W}_{\!K}\D{p_j}$):
\[
  (\W{W}_{\!V}\D P)(\W{W}_{\!K}\D P)^{\top}
   = \sum_{j=1}^{N}(\W{W}_{\!V}\D{p_j})(\W{W}_{\!K}\D{p_j})^{\top}.
\]

So attention over the whole prompt is a sum, over every prompt token
$\D{p_j}$, of (its value)\,$\otimes$\,(its key):
\[
  F(\D q)=\sum_{j=1}^{N}(\W{W}_{\!V}\D{p_j})(\W{W}_{\!K}\D{p_j})^{\top}\,\D q .
\]
Still line 1 --- just one outer-product term per token.

\section*{Step 2 — split the prompt's columns into base + demos}

The prompt $\D P$ is just columns. Cut them into two groups:
\[
  \D P=\big[\;\underbrace{\D{p_1}\ \cdots\ \D{p_{N_0}}}_{\text{base / instructions}}
        \;\;\big\|\;\;
        \underbrace{\D{x_1}\ \cdots\ \D{x_m}}_{\text{in-context examples}}\;\big]
\]

The sum over \emph{all} tokens splits into base tokens plus demo tokens:
\[
  F(\D q)=
  \underbrace{\sum_{j\in\text{base}}(\W{W}_{\!V}\D{p_j})(\W{W}_{\!K}\D{p_j})^{\top}}_{\textstyle \text{call this } \W{W}_{\mathrm{ZS}}}\,\D q
  \;+\;
  \sum_{i\in\text{demos}}(\W{W}_{\!V}\D{x_i})(\W{W}_{\!K}\D{x_i})^{\top}\,\D q .
\]

\section*{Step 3 — name the base part $\W{W}_{\mathrm{ZS}}$}

The first sum doesn't involve the examples at all --- it's the fixed matrix you
would get from the base prompt \emph{alone} (no demonstrations). That is the
``zero-shot'' model, so name it $\W{W}_{\mathrm{ZS}}$:
\[
  \W{W}_{\mathrm{ZS}}=\sum_{j\in\text{base}}(\W{W}_{\!V}\D{p_j})(\W{W}_{\!K}\D{p_j})^{\top}.
\]

Substitute the name, and you have line 2 exactly:
\[
  \boxed{\,F(\D q)=\W{W}_{\mathrm{ZS}}\,\D q+\sum_{i}(\W{W}_{\!V}\D{x_i})(\W{W}_{\!K}\D{x_i})^{\top}\,\D q\,}
\]

\hrulefill

\section*{The whole connection in one sentence}

Attention over the prompt is a sum of one (value\,$\otimes$\,key) term per
token (the matrix identity); split those tokens into base + examples, and the
base terms collapse into a fixed matrix $\W{W}_{\mathrm{ZS}}$ while the example
terms stay as an explicit sum --- which is line 2.

That leftover example-sum,
$\sum_i(\W{W}_{\!V}\D{x_i})(\W{W}_{\!K}\D{x_i})^{\top}$, is the patch
$\W{\Delta W}$ --- a weight matrix with the same outer-product shape as a
gradient step.

\medskip

\textbf{Optional slide fix.} Add the per-token line
$F(\D q)=\sum_j (\W{W}_{\!V}\D{p_j})(\W{W}_{\!K}\D{p_j})^{\top}\,\D q$
\emph{between} the current lines 1 and 2, so the split is shown rather than
implied.

\end{document}
